% Stage21 REVIEW ONLY: exact global rho=1 phase theorem and
% local endogenous-product-design missing-middle extension.
% Requires amsmath, amssymb, amsthm.
\section{An analytic missing-middle theorem for standards coalitions}
\label{sec:missing-middle}

Let $s$, $m(g)$, $o(g)$ and $i(g)$ denote the endogenous
continuation national payoffs for fragmentation, member of the
bilateral agreement, outsider to the bilateral agreement,
and grand standards respectively. Weak accession votes are
accepted on equality. In the inherited-depth extensive-form
game, a grand continuation is acceptable at $g$ if
$i(g)\ge\max\{s,m(g),o(g)\}$, whereas terminal SU requires
$m(g)\ge s$ and failure of at least one expansion vote.
This classification anticipates later accession in country B's
initial decision.

\begin{proposition}[No optimal bilateral union, and a network threshold]
\label{prop:no-SU-face}
Maintain the original three-country Cournot game, with $d>0$,
$\kappa>0$, no ex-post architecture adjustment friction
($\lambda=0$), $\rho=1$ and $0<\eta<1$. Then the sole
government-optimal terminal regime is
\[
 {\rm SW}\quad\text{if }0<\eta\le6/7,\qquad
 {\rm IS}\quad\text{if }6/7<\eta<1.
\]
SU is never proposer-optimal. For $\eta>6/7$ there exists
an explicitly constructed strictly acceptable IS offer
whose proposer's payoff EXCEEDS that under every
possible terminal SU, at every alternative depth.
\end{proposition}
\begin{proof}
When $\rho=1$, $\beta=1-\rho=0$. Every firm minimizes its
once-paid design cost at $x_i=1$ independent of government
history. Put
\[
 a=2-\eta,\quad
 \theta=\frac{1-\eta}{2-\eta}\in(0,1/2),\quad
 D=\frac{d^2}{a^2}>0.
\]
Solving the fulfilled Cournot equations gives
\[
 s=\frac92D,\quad
 m(g)=D\left\{\frac{4+g}{(1+\theta g)^2}+\frac12\right\},
\quad
 o(g)=D\left\{\frac{1+g}{(1+\theta g)^2}+\frac72\right\},
\quad
 i(g)=D\frac{\frac92+3g}{(1+2\theta g)^2}.
\]
In particular,
\[
 m(g)-s=
 D\frac{g(1-8\theta-4\theta^2g)}{(1+\theta g)^2}.
\]
If $\theta\ge1/8$ (equivalently $\eta\le6/7$), then
$m(g)<s$ at every $g>0$. Moreover,
\[
 i(g)-o(g)=-D\frac{gP(\theta,g)}
 {(1+\theta g)^2(1+2\theta g)^2},
\]
where
\[
 P=14\theta^4g^3+(42\theta^3+\theta^2)g^2+
 \theta(45\theta-2)g+(16\theta-2).
\]
Every term is nonnegative if $\theta\ge1/8$, and the
linear term is strictly positive for $g>0$, whence
$i(g)<o(g)$. SU is not initially acceptable, and
C blocks every IS accession: SW is the sole final regime.

Now suppose $0<\theta<1/8$ and write $v=1-8\theta>0$.
For every bilateral depth, $m(g)-s\le Dv$.
Choose the grand-offer depth
\[
 h=\min\{1,\,2v\}.
\]
We claim that $i(h)>s+Dv$, $i(h)>o(h)$ and $i(h)>m(h)$.
For $\theta\le1/16$, $h=1$. At that depth,
both $m(1)$ and $o(1)$ are bounded above by $11D/2$,
whereas $i(1)\ge160D/27>11D/2$ and
$i(1)-s\ge77D/54>Dv$.

For $1/16<\theta<1/8$, use $0<v<1/2$ and $h=2v$.
Set $A(v)=2(v-2)^2(v+1)^2$ and
$B(v)=A(v)(v^2-v-4)^2$, both positive. Exact algebra gives
\[
 \frac{i(h)-s-Dv}{D}
 =\frac{v(4+19v+24v^2-5v^3-2v^4)}{A(v)}>0,
\]
\[
 \frac{i(h)-o(h)}{D}
 =\frac{v^2(280+492v-83v^2-240v^3+
           42v^4+28v^5-7v^6)}{B(v)}>0,
\]
\[
 \frac{i(h)-m(h)}{D}
 =\frac{v(192+304v+132v^2-77v^3-
          48v^4+6v^5+4v^6-v^7)}{B(v)}>0.
\]
Indeed on $0<v<1/2$, the three collections of negative
higher-order monomials are bounded, respectively, by
$5v^3+2v^4<3/4<4$,
$83v^2+240v^3+7v^6<51<280$,
and $77v^3+48v^4+v^7<13<192$.
Thus this one accepted grand offer strictly outperforms
\emph{every} terminal bilateral outcome at any depth
and outperforms SW. The feasible IS set is nonempty and
closed in compact $[0,1]$ under weak-vote acceptance,
so an optimal IS depth is attained. It must prevail
over all SU and SW alternatives.
\end{proof}

\begin{proposition}[The missing middle persists under endogenous designs]
\label{prop:no-SU-endogenous}
Fix $d>0$, $\kappa>0$, $\lambda=0$ and any compact network
interval $K=[\eta_-,\eta_+]\Subset(0,1)$ with
$\eta_-<6/7<\eta_+$. There exist $\bar\rho<1$ and a
continuously differentiable threshold $\eta_c(\rho)$,
$\eta_c(1)=6/7$, such that for $\bar\rho<\rho<1$
and $\eta\in K$, the original firm's design Nash is unique,
interior and regular in every policy continuation and
\[
 \eta\le\eta_c(\rho)\implies {\rm SW},\qquad
 \eta>\eta_c(\rho)\implies {\rm IS}.
\]
In particular SU is never the government-optimal terminal
regime throughout this genuine endogenous-design parameter
strip.
\end{proposition}
\begin{proof}
At $\rho=1$ the full three-firm design pseudo-gradient is
$F_i=\kappa(1-x_i)$ and $F_x=-\kappa I_3$,
while all fulfilled Cournot quantities are positive with a
common strict nonsingularity margin over $K\times[0,1]$
and the three regimes. By compactness, a sufficiently small
$\beta=1-\rho>0$ preserves uniformly strict own concavity,
strong monotonicity, positive output, and a unique smooth
symmetric-by-permutation Nash continuation, now with
$0<x_i<1$; the unilateral design benefit has a strictly
inward gradient at the former boundary $x_i=1$.
Uniformly across the compact strategy cube, these conclusions
follow by continuity of the full Jacobian from $-\kappa I_3$.

All decisive sign inequalities in Proposition
\ref{prop:no-SU-face} persist for compact $\eta$ subsets
strictly above or below $6/7$: on the lower side, use
negative derivatives at $g=0$ as well as negative levels
for $g$ bounded away from zero; on the upper side, the
explicit grand-offer witness $h(\eta)$ has strict
accession and global SU-payoff-dominance margins.

It remains to cover the transition. Let $p=(\eta,\rho)$,
\[
 M(g,p)=m(g,p)-s(p),\quad
 V(g,p)=i(g,p)-o(g,p),\quad
 J(g,p)=i(g,p)-s(p),\quad
 H(g,p)=i(g,p)-m(g,p),
\]
and $\mu(p)=M_g(0,p)$. All four functions vanish
at $g=0$, and three-country edge symmetry supplies the
exact identity $V_g(0,p)=2\mu(p)$.
At $p_0=(6/7,1)$, $a=8/7$, $D=d^2/a^2$, one obtains
\[
 M(g,p_0)=-D\frac{g^2}{16(1+g/8)^2},\qquad
 V(g,p_0)=-D\frac{g^2(29/64+25g/256+7g^2/2048)}
 {(1+g/8)^2(1+g/4)^2},
\]
and $J_g(0,p_0)=H_g(0,p_0)=3D/4>0$.
Moreover
\[
 \mu(p_0)=0,\qquad
 \partial_\eta\mu(p_0)=\frac{8d^2}{(8/7)^4}
 =\frac{2401d^2}{512}>0.
\]
The implicit function theorem therefore defines a smooth
unique $\eta_c(\rho)$ locally with
$\mathrm{sign}\,\mu=\mathrm{sign}(\eta-\eta_c(\rho))$.

Uniform negativity of the quadratic terms at $p_0$,
continuity of the Hessians and compact negativity away
from $g=0$ yield constants $a_0,b_0,B_0,c_0,\epsilon>0$
on a common primitive neighborhood for which
\[
 M(g,p)\le\mu(p)g-a_0g^2,\quad
 2\mu(p)g-B_0g^2\le V(g,p)
 \le2\mu(p)g-b_0g^2,\quad
 J(g,p),H(g,p)\ge c_0g
\]
for $0\le g\le\epsilon$, while
$M(g,p),V(g,p)<0$ on $\epsilon\le g\le1$.
If $\mu\le0$, no positive-depth SU offer can win and
the outsider vetoes any IS accession: SW obtains.
If $\mu>0$, choose $h=\mu/B_0\le\epsilon$.
Then $V(h)\ge\mu^2/B_0>0$, $J(h),H(h)\ge c_0h>0$:
the IS proposal is strictly accepted. Across ALL bilateral
depths, $M(g,p)\le\mu^2/(4a_0)$; but
$J(h)\ge c_0\mu/B_0>\mu^2/(4a_0)$
after shrinking the primitive neighborhood.
Hence the accepted grand offer dominates every SU,
irrespective of its depth or actual veto status.
Weak-vote closedness and continuity give an attained
global IS optimal depth. An ordinary finite-cover argument
unifies these critical and away-from-critical neighborhoods
over $K$, completing the result.
\end{proof}

\paragraph{Scope.}
This is a rigorous two-regime theorem with strategic
endogenous firm differentiation for $\rho<1$ in a
nontrivial parameter strip. It is NOT the unrestricted
claim that SU can never be optimal for every
$\rho,\eta,\kappa,d$; nor does it characterize corner
Cournot equilibria outside the maintained model's
interior domain. The explicit grand depth $h$ is a
winning witness, not necessarily the maximizing depth.
Publication priority remains subject to close
proposition-level comparison with previous standards
coalition and network-effect models.
